\RequirePackage{iftex}
\ifPDFTeX
\documentclass[a4paper,12pt]{article}
\usepackage[margin=22mm]{geometry}
\usepackage[T1]{fontenc}
\usepackage[utf8]{inputenc}
\usepackage{graphicx}
\usepackage{CJKutf8}
\else
\documentclass[a4paper,12pt]{jsarticle}
\usepackage[dvipdfmx,margin=22mm]{geometry}
\usepackage[dvipdfmx]{graphicx}
\fi
\usepackage{amsmath,amssymb,amsthm}
\usepackage{enumitem}

\setlength{\parindent}{0pt}
\setlength{\parskip}{0.35em}
\setlength{\fboxsep}{8pt}
\setlist[itemize]{leftmargin=2em,itemsep=0.25em,topsep=0.35em}
\setlist[enumerate]{leftmargin=2.2em,itemsep=0.45em,topsep=0.35em}

\newcommand{\pointbox}[1]{%
\begin{center}
\fbox{\begin{minipage}{0.91\linewidth}#1\end{minipage}}
\end{center}
}

\begin{document}
\ifPDFTeX
\begin{CJK}{UTF8}{min}
\fi

\theoremstyle{definition}
\newtheorem*{definition}{定義}
\newtheorem*{theorem}{定理}
\newtheorem*{example}{例}
\newtheorem*{remark}{注意}
\renewcommand{\proofname}{\normalfont 証明}

\begin{center}
{\LARGE 数学1A 第11回レジュメ}\\[0.4em]
数列の極限と基本性質
\end{center}

\section*{1. 収束の定義}

\begin{definition}
数列 $\{a_n\}$ が $a$ に収束するとは、
\[
\forall \varepsilon>0\ \exists N\in\mathbb{N}\ \forall n>N\quad
|a_n-a|<\varepsilon
\]
が成り立つことをいう。
\end{definition}

\begin{remark}
既知の収束を使うときは、目標の $\varepsilon$ と区別して
$\varepsilon_1,\varepsilon_2$ などを用いると混乱しにくい。
\end{remark}

\section*{2. 和・積・商の極限}

\begin{theorem}
$a_n\to a,\ b_n\to b$ とする。このとき
\[
a_n+b_n\to a+b,\qquad a_nb_n\to ab
\]
が成り立つ。また $b\ne0$ ならば
\[
\frac{a_n}{b_n}\to \frac{a}{b}
\]
である。
\end{theorem}

\begin{proof}
和について示す。$\varepsilon>0$ を取る。$a_n\to a$ より、ある $N_1$ が存在して
$n>N_1$ ならば $|a_n-a|<\varepsilon/2$ である。同様に、ある $N_2$ が存在して
$n>N_2$ ならば $|b_n-b|<\varepsilon/2$ である。$N=\max\{N_1,N_2\}$ とすると
\[
|(a_n+b_n)-(a+b)|\le |a_n-a|+|b_n-b|<\varepsilon
\]
となる。

積については、まず $a_n\to a$ から十分大きい $n$ で $|a_n|\le |a|+1$ となる。
また
\[
|a_nb_n-ab|
\le |a_n|\,|b_n-b|+|b|\,|a_n-a|
\]
である。右辺を $\varepsilon$ より小さくするように
$|a_n-a|$ と $|b_n-b|$ を十分小さくすればよい。

商については、$b\ne0$ なので十分大きい $n$ で $|b_n|\ge |b|/2$ となる。この範囲で
\[
\left|\frac{a_n}{b_n}-\frac{a}{b}\right|
=\left|\frac{b(a_n-a)-a(b_n-b)}{bb_n}\right|
\]
を用いれば、積の場合と同様に示せる。
\end{proof}

\section*{3. 単調列}

\begin{definition}
$a_n\le a_{n+1}$ がすべての $n$ で成り立つとき、$\{a_n\}$ は単調増加である
という。$a_n\ge a_{n+1}$ がすべての $n$ で成り立つとき、単調減少であるという。
\end{definition}

\begin{theorem}[単調収束定理]
$\{a_n\}$ が単調増加で上に有界ならば、$\{a_n\}$ は収束し、
\[
\lim_{n\to\infty}a_n=\sup\{a_n\mid n\in\mathbb{N}\}
\]
である。
\end{theorem}

\begin{proof}
$\alpha=\sup\{a_n\mid n\in\mathbb{N}\}$ とおく。$\varepsilon>0$ を取る。
上限の近似性より、ある $N$ が存在して $\alpha-\varepsilon<a_N$ となる。
単調増加性より、$n>N$ ならば
\[
\alpha-\varepsilon<a_N\le a_n\le\alpha
\]
である。したがって $|a_n-\alpha|<\varepsilon$ であり、$a_n\to\alpha$ である。
\end{proof}

\section*{4. 部分列}

\begin{definition}
$n_1<n_2<\cdots$ を自然数の狭義単調増加列とする。このとき
\[
a_{n_1},a_{n_2},\ldots
\]
を $\{a_n\}$ の部分列という。
\end{definition}

\begin{theorem}
$a_n\to a$ ならば、任意の部分列 $a_{n_k}$ も $a$ に収束する。
\end{theorem}

\begin{proof}
$\varepsilon>0$ を取る。$a_n\to a$ より、ある $N$ が存在して $n>N$ ならば
$|a_n-a|<\varepsilon$ である。$n_k\ge k$ なので、$k>N$ ならば
$n_k>N$ である。したがって $|a_{n_k}-a|<\varepsilon$ となる。
\end{proof}

\begin{example}
$a_n=(-1)^n$ は収束しない。実際、偶数番目の部分列は $1$ に収束し、奇数番目の
部分列は $-1$ に収束する。もし元の数列が収束するなら、すべての部分列は同じ
極限に収束しなければならない。
\end{example}

\pointbox{
部分列は、収束しないことを示すときにも有効である。異なる極限をもつ部分列を
二つ見つければ、元の数列は収束しない。
}

\ifPDFTeX
\end{CJK}
\fi
\end{document}
